\subsection{When do divided comparisons help?}\label{sec:mechanism}
We first supplied the triplets directly so that nearest-neighbour search could not change their composition. Each world had six Gaussian channels measuring class, six measuring acquisition site and 48 measuring a shared background. Half the hit/miss assignments were reversed in expectation. Class and site favoured opposite references. When references shared the query background, noise divided the background votes. Otherwise a fair coin assigned the background match to one reference, giving a mean absolute row imbalance of 0.705 versus 0.101. The two regimes exchanged which comparisons had a shared background. In each of 48 independently generated worlds per regime, 192 acquired triplets selected six columns. We evaluated those columns on 768 fresh triplets and fitted logistic regression on 384 separate observations. In the 1,536 test observations class, site and background were independent.

When valid comparisons divided the background votes, $G_{16}$ improved mean test AUC over $G_0$ by 0.231 (97.5\% paired world-bootstrap interval 0.165 to 0.300). When reversed comparisons divided them, the difference was $-0.300$ ($-0.368$ to $-0.228$). Taking the minimum changed the $G_{16}$ support in 46 and 48 worlds, respectively. The assessment mean of the selected columns' weakest-layer scores \emph{fell} by 0.0054 and 0.0072. Test AUC changed by $-0.048$ and $+0.011$, the latter with a 95\% interval crossing zero. Here symmetry makes every population $G_0$ zero.

The instrument experiments used six stable and six polarity-reversing class channels with respective loadings of 0.55--0.85 and 1.30--1.70, plus 48 channels measuring one nuisance background (1.50--2.50). Gaussian measurement errors had unit standard deviation. Every hit label was correct. Queries used instrument A. With probability $1/3$ both references used A and shared the query background. Otherwise they used B, and only one matched the background. A logistic model fitted on 384 labelled observations (instrument B probability $2/3$, independent of class) served both instruments. With a fair background match, mean worse-instrument AUC across 48 worlds was 0.978 for the minimum, 0.075 for $G_{16}$ and 0.990 for $G_0$. A follow-up fixed four match probabilities before drawing another 48 worlds (Table~\ref{tab:acquisition}). Bias towards the hit raised nuisance-channel sign scores, while $G_{16}$ continued to favour reversing channels. At the $5/6$ policy, the minimum gained 0.254 over $G_{16}$ and 0.0118 over $G_0$ on the independent weakest-layer score (98.75\% intervals 0.245 to 0.264 and 0.0105 to 0.0131). Corresponding worse-instrument AUC gains were 0.881 and 0.408 (0.787 to 0.958 and 0.346 to 0.460). These are separate targets. The strongest single layer post hoc was $G_4$, 0.196 below the minimum. Using the labelled fitting sample, $F$ scored 0.991 and random selection 0.483.

\begin{table}[tb]
\centering\small
\caption{Mean of the lower of two instrument-specific test AUCs. Each column uses the same 48 new worlds and prediction samples. The policy is the probability that the hit matches the query's background in a cross-instrument comparison.}
\label{tab:acquisition}
\begin{tabular}{@{}lrrrr@{}}\toprule
Ranking & $1/2$ & $2/3$ & $5/6$ & $1$\\\midrule
$G_0$ & .990 & .947 & .567 & .489\\
$G_4$ & .887 & .856 & .779 & .650\\
$G_{16}$ & .094 & .094 & .094 & .094\\
Layer mean & .988 & .947 & .632 & .489\\
Minimum & .977 & .976 & .975 & .974\\\bottomrule
\end{tabular}
\end{table}
